\documentclass[border=5pt]{article}
\usepackage[margin=0pt,paperwidth=18cm,paperheight=8cm]{geometry}
\pagestyle{empty}
\usepackage{fontspec}
\usepackage{ctex}
\newcommand{\fontdir}{../../../00_知识库/论文写作/LaTeX模板_2026国赛版/fonts/}
\setCJKfamilyfont{song}[Path=\fontdir,UprightFont=SourceHanSerifCN-Regular.otf,BoldFont=SourceHanSerifCN-Bold.otf,AutoFakeBold=true]{SourceHanSerifCN}
\newcommand*{\song}{\CJKfamily{song}}
\usepackage{tikz}
\usetikzlibrary{arrows.meta,positioning,calc,patterns}
\definecolor{primary}{HTML}{1F3A93}\definecolor{primaryDark}{HTML}{1A3A6E}\definecolor{secondary}{HTML}{D9E4F2}
\definecolor{accent}{HTML}{D35400}\definecolor{accentLight}{HTML}{FDEBD0}
\definecolor{green}{HTML}{27AE60}\definecolor{red}{HTML}{C0392B}
\definecolor{gray}{HTML}{95A5A6}\definecolor{textDark}{HTML}{333333}\definecolor{arrowColor}{HTML}{5D6D7E}
\tikzset{
  arrow/.style={-{Stealth[length=2.0mm,width=1.8mm]},line width=0.8pt,draw=arrowColor},
  arrowR/.style={-{Stealth[length=2.0mm,width=1.8mm]},line width=1.0pt,draw=accent},
}

\begin{document}
\begin{center}
{\large\bfseries\song 图 7：Q1 求解原理示意图}\\[0.5cm]
\begin{tikzpicture}

  % ===== (a) 圆柱截面  =====
  \node[font=\song\small\bfseries] at (-5.5, 3.0) {(a) 圆柱截面与边界条件};

  \draw[primary,line width=1.2pt,fill=primary!4] (-5.5,0) circle (2.0cm);
  \fill[primary] (-5.5,0) circle (1.5pt);
  \draw[gray,line width=0.5pt] (-5.5,0) -- (-3.8,1.0);
  \node[font=\song\footnotesize] at (-3.5,1.4) {$R_0=2$ cm};
  \node[font=\song\footnotesize] at (-5.5,-0.4) {$\partial_r T=\partial_r C=0$};

  \draw[arrowR] (-3.4,0.3) -- (-2.2,1.2);
  \node[font=\song\footnotesize,text=accent,anchor=west] at (-2.2,1.2) {Robin 边界};

  \draw[arrow,draw=secondary!80!accent,line width=0.8pt] (-4.6,0.5) -- (-5.0,0.1);
  \draw[arrow,draw=secondary!80!accent,line width=0.8pt] (-5.1,0.0) -- (-5.3,-0.2);
  \node[font=\song\footnotesize,text=accent] at (-7.0,0.6) {热/湿流（外$\to$内）};

  % ===== (b) 控制体  =====
  \node[font=\song\small\bfseries] at (0.3, 3.0) {(b) 径向控制体网格};

  \draw[primary,line width=1.0pt] (0.3,0) circle (2.0cm);
  \foreach \r in {0.5,1.0,1.5} {
    \draw[gray,dashed,line width=0.3pt] (0.3,0) circle (\r cm);
  }
  \foreach \a in {0,45,90,135,180,225,270,315} {
    \draw[gray,dashed,line width=0.25pt] (0.3,0) -- ({0.3+2*cos(\a)},{2*sin(\a)});
  }

  \fill[green!15,draw=green,line width=0.7pt] (0.9,0.1) rectangle (1.3,0.3);
  \node[font=\song\footnotesize,text=green] at (2.1,1.1) {控制体};

  \fill[primary!12,draw=primary,line width=0.7pt] (0.3,0) circle (0.3cm);
  \node[font=\song\footnotesize,text=primary] at (-1.2,1.2) {中心控制体（$\Delta r/2$）};

  \node[font=\song\footnotesize] at (0.3,-2.6) {$N=80$ 区间，$\Delta r=0.25$ mm};

  % ===== (c) 时空离散  =====
  \node[font=\song\small\bfseries] at (6.5, 3.0) {(c) 时空离散与推进};

  % 网格
  \foreach \i in {0,1,2,3,4} {
    \draw[gray,dashed,line width=0.3pt] (4.0, \i*0.6-1.2) -- (8.5, \i*0.6-1.2);
  }
  \foreach \j in {0,1,2,3,4,5,6} {
    \draw[gray,dashed,line width=0.3pt] (4.0+\j*0.65, -1.2) -- (4.0+\j*0.65, 1.2);
  }

  % 节点
  \foreach \j in {0,1,2,3,4,5} {
    \foreach \i in {1,2,3} {
      \fill[primary] (4.3+\j*0.65, \i*0.6-1.2) circle (1.0pt);
    }
  }

  % 已知点
  \fill[accent] (7.5, 0.0) circle (2.0pt);
  \node[font=\song\footnotesize,text=accent] at (7.5, 0.3) {已知 $t^n$};

  % 推进方向
  \draw[arrowR,line width=1.0pt] (4.5, 1.6) -- (7.0, 1.6);
  \node[font=\song\footnotesize,text=accent] at (5.7, 1.9) {时间推进方向 $\to$};

  \node[font=\song\footnotesize] at (9.0, 0.3) {$\Delta t=1$ s（输出）};
  \node[font=\song\footnotesize] at (9.0, -0.3) {子步 $1/32$ s};

\end{tikzpicture}
\end{center}
\end{document}